\documentclass[11pt,a4paper]{article}

\usepackage[a4paper,margin=1in]{geometry}
\usepackage{fontspec}
\IfFontExistsTF{Times New Roman}
  {\setmainfont{Times New Roman}}
  {\IfFontExistsTF{TeX Gyre Termes}{\setmainfont{TeX Gyre Termes}}{\setmainfont{Times}}}
\usepackage{newtxmath}
\usepackage{amsmath}
\usepackage{booktabs}
\usepackage{longtable}
\usepackage{array}
\usepackage{enumitem}
\usepackage{xcolor}
\usepackage{listings}
\usepackage{graphicx}
\usepackage{float}
\usepackage[hidelinks]{hyperref}
\usepackage{xurl}
\usepackage{fancyhdr}
\usepackage{microtype}

\setlength{\parindent}{0pt}
\setlength{\parskip}{0.65em}
\setlength{\emergencystretch}{3em}
\setlist[itemize]{leftmargin=1.4em,itemsep=0.2em,topsep=0.2em}
\setlist[enumerate]{leftmargin=1.6em,itemsep=0.2em,topsep=0.2em}

\pagestyle{fancy}
\fancyhf{}
\lhead{\small Ticket 10A Study Report}
\rhead{\small Rannacher Smoothing Comparison Gate}
\cfoot{\thepage}
\renewcommand{\headrulewidth}{0.3pt}

\lstset{
  basicstyle=\ttfamily\small,
  breaklines=true,
  frame=single,
  columns=fullflexible,
  keepspaces=true,
  rulecolor=\color{black!30},
  xleftmargin=0.5em,
  xrightmargin=0.5em
}

\newcommand{\code}[1]{\texttt{#1}}
\newcommand{\file}[1]{\path{#1}}
\newcommand{\func}[1]{\path{#1}}
\newcolumntype{L}[1]{>{\raggedright\arraybackslash}p{#1}}
\newcolumntype{P}[1]{>{\raggedright\arraybackslash\ttfamily\footnotesize}p{#1}}

\title{\textbf{Ticket 10A Study Report: Rannacher Smoothing Comparison Gate}\\
\large SURF2026 American Option Risk Surfaces}
\author{Codex-assisted implementation report}
\date{June 16, 2026}

\begin{document}
\maketitle

\begin{abstract}
Ticket 10A adds a controlled Rannacher smoothing comparison after the Ticket 10 Delta and
Gamma diagnostics. It does not replace the validated Ticket 04 American CN/PSOR baseline.
Instead, it creates a named Rannacher variant that uses two backward-Euler half-substeps over
the first ordinary time interval, then resumes ordinary CN/PSOR. The experiment compares
baseline and smoothed results for price levels, PSOR convergence, threshold-based boundary
curves, and finite-difference Gamma diagnostics. The gate conclusion for the representative
American put and dividend-paying American call cases is: keep the baseline for now. The
smoothed variant changed prices and boundaries only slightly, but it did not improve the
headline strict-mask Gamma metrics in this run.
\end{abstract}

\tableofcontents
\newpage

\section{Role of Ticket 10A in the Whole Solver Project}

Tickets 01 through 10 built a validation ladder before this comparison gate:
\begin{itemize}
  \item Ticket 01 created payoff and Black--Scholes closed-form utilities.
  \item Tickets 02 and 03 built the grid, spatial operator, and European CN validation mode.
  \item Ticket 04 implemented the American CN/PSOR obstacle core.
  \item Tickets 05 through 07 validated American put, no-dividend American call, and dividend-paying American call behavior.
  \item Ticket 08 added obstacle and complementarity diagnostics.
  \item Ticket 09 extracted approximate free-boundary curves from \(U-\Phi\).
  \item Ticket 10 computed finite-difference Delta and Gamma diagnostics and warned that Gamma is fragile near payoff kinks, maturity, and extracted boundary regions.
\end{itemize}

Ticket 10A asks whether a Rannacher-style startup damping step improves those fragile Greek
diagnostics enough to justify using it later. It is a comparison gate, not a new default solver.
The baseline solver in \file{src/american_risk_surfaces/solvers/cn_psor.py} remains unchanged.

\section{Theory and Methodology Connection}

The documented source basis for this ticket is internal:
\begin{itemize}
  \item the literature map identifies Rannacher smoothing as a classical finite-difference stabilization idea near payoff kinks;
  \item the formulation note explains American option values, payoff obstacles, and Greeks;
  \item the solver validation plan requires an explicit unsmoothed-versus-smoothed comparison before smoothing is adopted;
  \item Ticket 09 and Ticket 10 provide boundary and Greek diagnostic tools used by this comparison.
\end{itemize}

No BibTeX file is created. The report does not invent external bibliographic details. It uses
the project documents as the source basis and treats Rannacher smoothing as a numerical
time-stepping device, not an economic model change.

\section{Rannacher Smoothing in Plain Language}

Crank--Nicolson is accurate and stable for smooth solutions, but option payoffs are not smooth at
the strike. At maturity the value equals payoff, so the first time steps start from a kinked
function. Gamma is a second derivative, so it can amplify that kink into noisy or exaggerated
curvature.

Rannacher smoothing handles this by beginning with a few backward-Euler substeps. Backward Euler
is more damping than Crank--Nicolson. The idea is:
\begin{enumerate}
  \item begin at maturity payoff;
  \item take a small number of implicit backward-Euler steps to damp the kink;
  \item continue with ordinary Crank--Nicolson after the startup interval.
\end{enumerate}

This should not be used just to make a plot look smoother. It must be checked against prices,
boundaries, and solver convergence.

\section{Why This Comes After Greek Diagnostics}

Ticket 10 first measured the baseline Delta and Gamma behavior. That made the problem visible:
Gamma can be fragile near maturity, payoff kinks, and extracted boundaries. Ticket 10A then tests
one targeted numerical remedy. This order matters. Without Ticket 10's diagnostics, smoothing
could become a vague aesthetic preference. With Ticket 10's diagnostics, smoothing can be judged
using evidence.

\section{Backward-Euler Substeps Versus Crank--Nicolson}

The ordinary American CN/PSOR step solves a linear complementarity problem based on
\[
A_{\mathrm{CN}} = I - \frac{1}{2}\Delta\tau L_h,
\]
with a right-hand side containing
\[
\left(I + \frac{1}{2}\Delta\tau L_h\right)U^n.
\]

The Ticket 10A startup substep instead uses backward Euler:
\[
A_{\mathrm{BE}} = I - \delta\tau L_h.
\]
For each substep, the interior LCP is
\[
A_{\mathrm{BE}}u \ge b,\qquad u\ge \Phi,\qquad
(u-\Phi)^\top(A_{\mathrm{BE}}u-b)=0,
\]
where \(b\) is the previous value vector adjusted for the new-time boundary values. Ticket 10A's
default is two equal substeps over the first ordinary interval:
\[
\delta\tau=\frac{\Delta\tau}{2}.
\]

\section{Implementation Design Decisions}

\begin{longtable}{P{0.32\textwidth}L{0.50\textwidth}}
\toprule
Decision & Reason \\
\midrule
\endhead
\path{rannacher.py} is separate & Prevents silent changes to the validated Ticket 04 baseline solver. \\
Two half-substeps & Matches the planned Rannacher gate and keeps the comparison simple. \\
\path{rannacher_substeps=0} delegates to baseline & Provides an explicit off switch and a testable equivalence path. \\
Reuse Ticket 04 PSOR & Keeps projection logic identical to the American obstacle core. \\
Reuse Ticket 09 and 10 diagnostics & Boundary and Gamma comparisons use existing validated diagnostic layers. \\
Gate recommendation is conservative & Smoothing is adopted only if it improves Greeks without material price or boundary distortion. \\
\bottomrule
\end{longtable}

\section{Code-Level Function Explanations}

\begin{longtable}{P{0.36\textwidth}L{0.46\textwidth}}
\toprule
Function or class & Purpose \\
\midrule
\endhead
\func{RannacherSmoothingMetadata} & Records whether smoothing is enabled, number and size of substeps, separated PSOR metadata, convergence status, and obstacle violation. \\
\func{RannacherCNPSORResult} & Wraps a standard \func{AmericanCNPSORResult} with Rannacher-specific metadata. \\
\func{backward_euler_psor_step} & Solves one implicit-Euler American LCP using the existing PSOR projection solver. \\
\func{rannacher_crank_nicolson_psor_price} & Runs two backward-Euler half-substeps, stores the first main time row, then continues with ordinary CN/PSOR. \\
\file{experiments/08_rannacher_smoothing_comparison.py} & Runs baseline and smoothed representative cases, writes CSV comparisons, and creates figures. \\
\bottomrule
\end{longtable}

\section{Comparison Cases and Parameters}

\begin{longtable}{P{0.25\textwidth}P{0.12\textwidth}L{0.43\textwidth}}
\toprule
Case & Type & Parameters \\
\midrule
\endhead
\path{american_put_medium} & put & \(K=1.0\), \(T=1.0\), \(r=0.05\), \(q=0.02\), \(\sigma=0.20\), \(S_{\max}=4.0\), \(M=N=120\). \\
\path{dividend_call_medium} & call & \(K=1.0\), \(T=1.0\), \(r=0.05\), \(q=0.08\), \(\sigma=0.20\), \(S_{\max}=4.0\), \(M=N=120\). \\
\bottomrule
\end{longtable}

\begin{longtable}{P{0.32\textwidth}L{0.46\textwidth}}
\toprule
Scheme setting & Value \\
\midrule
\endhead
Ordinary time interval & \(\Delta\tau=1/120=0.00833333333333\). \\
Rannacher substeps & \(2\). \\
Substep size & \(0.00416666666667\). \\
Post-startup method & Ordinary CN/PSOR for the remaining main time intervals. \\
Baseline behavior & Unchanged Ticket 04 American CN/PSOR solver. \\
\bottomrule
\end{longtable}

\section{Price Comparison Tables and Interpretation}

\begin{longtable}{P{0.24\textwidth}P{0.14\textwidth}P{0.14\textwidth}P{0.14\textwidth}P{0.16\textwidth}}
\toprule
Case & Max abs price diff & RMSE diff & Max-diff spot & Gate \\
\midrule
\endhead
\path{american_put_medium} & \(1.66937904909\times 10^{-6}\) & \(4.19362587648\times 10^{-7}\) & \(1.000000\) & \path{KEEP_BASELINE_FOR_NOW} \\
\path{dividend_call_medium} & \(1.72323866424\times 10^{-6}\) & \(4.26128419836\times 10^{-7}\) & \(1.033333\) & \path{KEEP_BASELINE_FOR_NOW} \\
\bottomrule
\end{longtable}

The price differences are very small in both representative cases. This means the Rannacher
startup does not materially distort the price level in this run. However, price similarity alone
does not justify adopting smoothing. The Greek and boundary comparisons must also support the
change.

\begin{longtable}{P{0.24\textwidth}P{0.12\textwidth}P{0.12\textwidth}P{0.14\textwidth}P{0.14\textwidth}P{0.14\textwidth}}
\toprule
Case & \(S/K\) & Payoff & Baseline & Rannacher & Difference \\
\midrule
\endhead
\path{american_put_medium} & 0.8 & 0.200000 & 0.200381130316 & 0.200380822124 & \(-3.0819\times10^{-7}\) \\
\path{american_put_medium} & 1.0 & 0.000000 & 0.066301330539 & 0.066299661160 & \(-1.6694\times10^{-6}\) \\
\path{american_put_medium} & 1.2 & 0.000000 & 0.015983679585 & 0.015982898743 & \(-7.8084\times10^{-7}\) \\
\path{dividend_call_medium} & 0.8 & 0.000000 & 0.008367161513 & 0.008366693230 & \(-4.6828\times10^{-7}\) \\
\path{dividend_call_medium} & 1.0 & 0.000000 & 0.065131176799 & 0.065129465900 & \(-1.7109\times10^{-6}\) \\
\path{dividend_call_medium} & 1.2 & 0.200000 & 0.203695017872 & 0.203694325038 & \(-6.9283\times10^{-7}\) \\
\bottomrule
\end{longtable}

\section{Boundary Comparison Tables and Interpretation}

\begin{longtable}{P{0.28\textwidth}P{0.13\textwidth}P{0.13\textwidth}P{0.13\textwidth}P{0.14\textwidth}}
\toprule
Case & Matched rows & Max shift & Mean shift & Status \\
\midrule
\endhead
\path{american_put_medium} & 120 & 0.000540799953 & \(1.09039294769\times10^{-5}\) & \path{COMPARABLE} \\
\path{dividend_call_medium} & 120 & 0.000663055353 & \(1.48912617221\times10^{-5}\) & \path{COMPARABLE} \\
\bottomrule
\end{longtable}

The threshold-based boundary curves remain close. This is reassuring, but it is not proof of
exact free-boundary accuracy. Boundary extraction inherits the Ticket 09 threshold and grid
uncertainty. The important gate point is that smoothing did not create an obvious boundary
distortion in these cases.

\section{Greek Comparison Tables and Interpretation}

\begin{longtable}{P{0.27\textwidth}P{0.16\textwidth}P{0.16\textwidth}P{0.16\textwidth}P{0.13\textwidth}}
\toprule
Case & Baseline full max \(|\Gamma|\) & Rannacher full max \(|\Gamma|\) & Baseline strict max \(|\Gamma|\) & Strict change \\
\midrule
\endhead
\path{american_put_medium} & 30.000000 & 30.000000 & 3.180018866 & \(+0.000266535\) \\
\path{dividend_call_medium} & 30.000000 & 30.000000 & 2.679887372 & \(+0.000341023\) \\
\bottomrule
\end{longtable}

The full-grid maximum \(|\Gamma|\) is unchanged because the strongest values remain at fragile
regions such as maturity or kink bands. The strict-mask Gamma metric is slightly higher after
smoothing in both cases. This does not mean Rannacher smoothing is bad in general. It means this
particular comparison does not provide enough evidence to adopt it as the default or even as the
preferred Greek-diagnostic path.

\section{PSOR Convergence Comparison}

Both baseline and Rannacher runs converged for both representative cases. The Rannacher variant
records the startup backward-Euler PSOR solves separately from the later CN/PSOR solves. This is
important because a smoothed variant would be risky if it improved Gamma while causing hidden
non-convergence. That did not occur here.

\begin{longtable}{P{0.28\textwidth}P{0.18\textwidth}P{0.18\textwidth}P{0.18\textwidth}}
\toprule
Case & Baseline converged & Rannacher converged & Obstacle violation \\
\midrule
\endhead
\path{american_put_medium} & True & True & 0 for both methods \\
\path{dividend_call_medium} & True & True & 0 for both methods \\
\bottomrule
\end{longtable}

\section{Figure Interpretation}

\begin{figure}[H]
  \centering
  \includegraphics[width=0.92\textwidth]{results/01_solver_validation/figures/ticket_10a_value_profiles.png}
  \caption{Baseline and Rannacher final value profiles nearly overlap for both cases.}
\end{figure}

\begin{figure}[H]
  \centering
  \includegraphics[width=0.92\textwidth]{results/01_solver_validation/figures/ticket_10a_price_difference_profiles.png}
  \caption{Smoothed-minus-baseline price differences are very small across moneyness.}
\end{figure}

\begin{figure}[H]
  \centering
  \includegraphics[width=0.92\textwidth]{results/01_solver_validation/figures/ticket_10a_gamma_profiles.png}
  \caption{Gamma profiles show why the comparison is a diagnostic gate rather than a visual smoothing exercise.}
\end{figure}

\begin{figure}[H]
  \centering
  \includegraphics[width=0.92\textwidth]{results/01_solver_validation/figures/ticket_10a_gamma_full_vs_strict.png}
  \caption{Full-grid and strict-mask max \(|\Gamma|\) over time for baseline and Rannacher runs.}
\end{figure}

\begin{figure}[H]
  \centering
  \includegraphics[width=0.92\textwidth]{results/01_solver_validation/figures/ticket_10a_boundary_curve_comparison.png}
  \caption{Threshold-based boundary curves are close, but remain approximate diagnostics.}
\end{figure}

\begin{figure}[H]
  \centering
  \includegraphics[width=0.92\textwidth]{results/01_solver_validation/figures/ticket_10a_psor_iterations.png}
  \caption{PSOR iteration counts remain well behaved for baseline and Rannacher solves.}
\end{figure}

\section{Testing Strategy}

\begin{longtable}{P{0.34\textwidth}L{0.47\textwidth}}
\toprule
Test category & Interpretation \\
\midrule
\endhead
Backward-Euler PSOR step & Checks shape, obstacle projection, boundary placement, and convergence metadata. \\
Rannacher result shape & Verifies that the smoothed result preserves the baseline grid and value-grid shape. \\
Zero substeps & Confirms that \path{rannacher_substeps=0} delegates to baseline behavior. \\
Obstacle checks & Confirms smoothed put and dividend-call values remain above payoff. \\
Price closeness & Checks that smoothed and baseline prices remain finite and close on representative cases. \\
Boundary and Greek diagnostics & Verifies that Ticket 09 and Ticket 10 diagnostics run on both result types. \\
Artifact checks & Confirms CSV columns and figure files are created. \\
Scope checks & Rejects stress-map, dataset, neural-network, or machine-learning-label scope creep. \\
\bottomrule
\end{longtable}

\section{Test Results and Interpretation}

The final integrated command was:
\begin{lstlisting}
PYTHONPYCACHEPREFIX=/private/tmp/codex_pycache PYTHONPATH=src MPLCONFIGDIR=/private/tmp/matplotlib-cache FC_CACHEDIR=/private/tmp/fontconfig-cache python3 -m unittest discover -s tests
\end{lstlisting}

Result:
\begin{lstlisting}
Ran 135 tests in 2.901s
OK
\end{lstlisting}

The Ticket 10A-focused test file also passed:
\begin{lstlisting}
Ran 14 tests in 0.895s
OK
\end{lstlisting}

These tests show that the named Rannacher variant, comparison experiment, CSV artifacts, and
figures are reproducible in the current environment. They do not prove that Rannacher smoothing
is globally better. That would require grid/domain sensitivity and broader parameter testing.

\section{Implementation Pitfalls and Debugging Notes}

Actual recorded implementation iterations:
\begin{longtable}{P{0.24\textwidth}L{0.24\textwidth}L{0.24\textwidth}L{0.17\textwidth}}
\toprule
Event & Symptom & Fix & Beginner lesson \\
\midrule
\endhead
Initial TDD red run & Full suite reported 9 errors because \path{rannacher.py} and the Ticket 10A experiment did not exist yet. & Implemented the new module and experiment files. & Red tests are useful when they fail for the expected missing capability. \\
Focused module invocation & \path{python -m unittest tests.test_rannacher_smoothing...} failed because \path{tests/} is not a package. & Used unittest discovery with \path{-s tests -p test_rannacher_smoothing.py}. & Test invocation style must match repository structure. \\
\bottomrule
\end{longtable}

Likely traps to watch for in future Rannacher work:
\begin{longtable}{P{0.34\textwidth}L{0.46\textwidth}}
\toprule
Pitfall & Why it matters \\
\midrule
\endhead
Silently changing the baseline solver & It would invalidate earlier validation tickets and hide the comparison. \\
Claiming smoothing is better because plots look smoother & The gate must check prices, boundaries, Greeks, and convergence. \\
Ignoring price shifts & A smoother Gamma is not acceptable if prices move materially. \\
Ignoring boundary shifts & Boundary curves are approximate, but large shifts are still warning signs. \\
Comparing different grids or parameters & The comparison would mix smoothing effects with discretization effects. \\
Using smoothed Greeks as dataset labels too early & Dataset and neural-surface work is out of scope and requires later validation. \\
\bottomrule
\end{longtable}

\section{Gate Conclusion with Evidence}

The gate conclusion for Ticket 10A is:
\[
\boxed{\text{\code{KEEP\_BASELINE\_FOR\_NOW}}}
\]

Evidence:
\begin{itemize}
  \item price differences are very small, around \(1.7\times10^{-6}\);
  \item boundary shifts are small, below \(7\times10^{-4}\);
  \item PSOR convergence remains clean;
  \item strict-mask Gamma does not improve in either representative case.
\end{itemize}

Therefore, the Rannacher variant is useful as a named comparison tool, but the project should not
adopt it as the default or claim that it improves Greek diagnostics yet. A later grid/domain
sensitivity ticket can revisit this decision with a broader evidence base.

\section{Limitations}

\begin{longtable}{P{0.32\textwidth}L{0.48\textwidth}}
\toprule
Limitation & Explanation \\
\midrule
\endhead
Only two representative cases & The comparison is not full grid/domain sensitivity. \\
Only one smoothing setting & The default uses two half-substeps; other choices are not explored. \\
Boundary curves are threshold-based & Boundary shifts inherit Ticket 09 extraction uncertainty. \\
Greek masks are diagnostic & Strict-mask Gamma is useful, but not an analytical proof. \\
No stress maps or datasets & Ticket 10A intentionally avoids later machine-learning or stress-surface work. \\
\bottomrule
\end{longtable}

\section{Lessons Learned / Study Notes}

For beginner researchers, the main lesson is that numerical smoothing is a trade-off. It may damp
payoff-kink artifacts, but it is not automatically better. A careful comparison asks:
\begin{itemize}
  \item Did prices stay close?
  \item Did obstacle and PSOR convergence remain clean?
  \item Did approximate boundaries shift?
  \item Did Gamma actually improve under the chosen diagnostic mask?
\end{itemize}

In this ticket, prices and boundaries stayed close, but Gamma did not improve. That is a valid
and useful result.

\section{Link to the Next Ticket}

The next ticket should continue the validation ladder without treating Ticket 10A as a default
solver change. A natural next step is grid/domain sensitivity: compare solver outputs across
mesh sizes and domain cutoffs before using Greeks, boundaries, or prices as downstream research
features.

\section{Reproducibility Record}

\begin{longtable}{P{0.35\textwidth}L{0.45\textwidth}}
\toprule
Item & Record \\
\midrule
\endhead
New source file & \file{src/american_risk_surfaces/solvers/rannacher.py}. \\
New experiment & \file{experiments/08_rannacher_smoothing_comparison.py}. \\
New test file & \file{tests/test_rannacher_smoothing.py}. \\
CSV outputs & Four files under \file{results/01_solver_validation/tables/} with prefix \path{ticket_10a}. \\
Figure outputs & Six files under \file{results/01_solver_validation/figures/} with prefix \path{ticket_10a}. \\
Report outputs & This \LaTeX{} source and compiled PDF under \file{reports/03_solver/tickets/}. \\
\bottomrule
\end{longtable}

\end{document}
